\documentclass{article}
\usepackage{amsmath,amssymb}
\title{Positive upper density of restricted sumsets: Erd\H{o}s 339, second question}
\author{Formalization extension developed with an AI assistant}
\date{}
\begin{document}
\maketitle
\paragraph{Attribution.} The mathematical result is due to N. Hegyv\'ari,
F. Hennecart and A. Plagne, J. Reine Angew. Math. 560 (2003), 199--220.
The finite estimate below is reused from the formalization by Codex and
GPT-5.6 Sol in plby/lean-proofs, revision
\texttt{8822f7ddef30fadbd92e1c6ab4ed897af356af5e}.
This contribution supplies the upper-density deduction, not a new mathematical discovery.

\paragraph{Statement.} For $A\subseteq\mathbb N$ and $r\ge1$, let $rA$
be the set of sums of $r$ members of $A$, allowing repetition, and let
$r^{\wedge}A$ require the $r$ summands to be pairwise distinct.
Write $S(N)=|S\cap[0,N)|$ and
$\overline d(S)=\limsup_{N\to\infty}S(N)/N$.
Then $\overline d(rA)>0$ implies $\overline d(r^{\wedge}A)>0$.

\paragraph{Finite estimate.} Put $q=2^r$ and
$K=2^{r^2}(1+r(r+1))^r$.
The upstream finite equality-pattern argument proves
\[
 (rA)(N)\le K\,(r^{\wedge}A)(qN)
 \qquad\text{whenever }A(N)\ge r+1.
\]
For completeness, the counting mechanism is as follows. Given bases $b$
and a reservoir $X$ of at least $r+1$ possible new summands, forbid at most
$r$ partners for each base. Choose $r+1$ reservoir points. Every base allows
at least one of them; translation by each fixed point is injective.
Thus the number of bases is at most $(r+1)$ times the number of allowed sums.
At most $r$ forbidden sums per base are lost, giving the factor $1+r(r+1)$.
Partition representing tuples by their equality matrices (at most $2^{r^2}$).
Apply the forbidden-partner bound to isolate each coordinate in turn;
each step doubles the cutoff, and after $r$ steps all coordinates are distinct.

\paragraph{Upper-density deduction.} If $\overline d(rA)>0$, then $A$
is infinite, since a finite $A$ has finite $rA$ and hence zero upper density.
Consequently $A(N)\ge r+1$ eventually. Dividing the finite estimate by $N>0$ gives
\[
 \frac{(rA)(N)}N\le Kq\frac{(r^{\wedge}A)(qN)}{qN}.
\]
The limsup along the cofinal subsequence $qN$ is no greater than the full
limsup. Therefore
\[
 \overline d(rA)\le Kq\,\overline d(r^{\wedge}A).
\]
Since $Kq>0$, the conclusion follows. The formal endpoint also permits
$r=0$: its premise is then impossible, because $0A=\{0\}$ has density zero.

\paragraph{Scope.} This is the original second question, not merely the
first question about asymptotic bases and lower density. No claim of novelty,
independent verification, or award eligibility is made.
\end{document}
